% ===========================================================================
% ACT VI -- Tooling, the GPU Integrator, and the Dataset
% (Ch. 6 + a chapter that does not exist yet)   Slides 49-62
% Budget: 8 minutes. The strongest completed quantitative results in the deck.
%
% Slide 55 was CORRECTED 2026-09-22: the earlier version claimed the torch port
% was a dependency of the model. It is not. See NOTES.md.
% ===========================================================================

\actdivider{6}{Tooling, GPU, and Data}

% --- 49. Why a console ----------------------------------------------------
\begin{frame}{Why a Graphical Console}
  \vspace{0.6em}
  \centering
  \scalebox{0.70}{\raillabLayers}

  \vspace{1.2em}
  \begin{minipage}{0.90\textwidth}\small
    Building a scenario means checking the train and its control commands
    \textbf{visually} before trusting them, which is easier in a graphical tool
    than in scripts. One design rule keeps the tool separate: \alert{only the
    service layer calls the simulation code}. The simulator is unchanged by the
    application and still runs without it.
  \end{minipage}
\end{frame}

% --- 50. Consist screen ---------------------------------------------------
\begin{frame}{RailLab: Building a Consist}
  \vspace{0.2em}
  \begin{columns}[c,onlytextwidth]
    \begin{column}{0.63\textwidth}
      \centering
      \includegraphics[width=\linewidth,height=0.74\textheight,keepaspectratio]%
        {../thesis_imgs_7_15/raillab/raillab_consist.png}
    \end{column}
    \begin{column}{0.34\textwidth}
      \small
      A library of reusable car types, the ordered train under construction with
      per-car overrides, and saved trains.

      \vspace{0.8em}
      The running summary reports car count, total mass, total length and
      powered vehicles.

      \vspace{0.8em}
      Backing it: \textbf{10 car types} and roughly \textbf{500 saved
      consists}.
    \end{column}
  \end{columns}
\end{frame}

% --- 51. Controls screen --------------------------------------------------
\begin{frame}{RailLab: Control Profiles and Distributed Power}
  \vspace{0.2em}
  \begin{columns}[c,onlytextwidth]
    \begin{column}{0.56\textwidth}
      \centering
      \includegraphics[width=\linewidth,height=0.76\textheight,keepaspectratio]%
        {../thesis_imgs_7_15/raillab/raillab_controls.png}
    \end{column}
    \begin{column}{0.41\textwidth}
      \small
      Randomization bounds and distributed-power settings, with a live preview
      of the traction command per power role and the shared brake command.

      \vspace{0.9em}
      \alert{The offset between the three traction traces} is the configured lag
      of the mid-train and rear helper units behind the head end.
    \end{column}
  \end{columns}
\end{frame}

% --- 52. Results screen ---------------------------------------------------
\begin{frame}{RailLab: A 126-Vehicle Consist, 120 Seconds}
  \vspace{0.1em}
  \centering
  \includegraphics[width=0.74\linewidth,height=0.66\textheight,keepaspectratio]%
    {../thesis_imgs_7_15/raillab/sim_results.png}

  \vspace{0.3em}
  \begin{minipage}{0.92\textwidth}\footnotesize
    Every vehicle's speed above every coupler's force, on the same time axis.
    \textbf{The spread between vehicle speeds} widens and narrows as slack runs in
    and out, and the coupler forces below spread out at the same times.
  \end{minipage}
\end{frame}

% --- 53. Driving regimes --------------------------------------------------
\begin{frame}{Driving Regimes Instead of Random Commands}
  \vspace{0.1em}
  \centering
  \scalebox{0.75}{\regimeGrid}

  \vspace{0.1em}
  \begin{minipage}{0.94\textwidth}\scriptsize
    Fully random commands mostly describe driving no operator would do. Instead,
    commands come from \textbf{11 typical manoeuvres}. A check after generation
    ensures traction and brake are not applied together, with
    \alert{0 violations across 5,500 profiles}. Two exceptions are intended:
    \emph{stretch braking} applies brake against traction, and in an
    \emph{emergency} the two overlap while traction fades.
  \end{minipage}
\end{frame}

% --- 54. The NumPy simulator could not build this ------------------------
\begin{frame}{The NumPy Simulator Could Not Build This Dataset}
  \vspace{-0.2em}
  \begin{columns}[T,onlytextwidth]
    \begin{column}{0.58\textwidth}
      \benchRegimeChart
      \vspace{-0.2em}
      \centering\benchLegend
    \end{column}
    \begin{column}{0.40\textwidth}
      \vspace{0.2em}
      {\footnotesize Two causes, both fixed by moving the simulator to the GPU:}
      \vspace{0.4em}

      \begin{itemize}\setlength\itemsep{0.45em}
        \item[] {\footnotesize\textbf{1.} The equations used a
              \textbf{Python loop} over vehicles: 130 slow iterations per
              evaluation at $N = 130$.}
        \item[] {\footnotesize\textbf{2.} Force changes abruptly at the edge of
              the slack, so the solver shrank its step whenever \emph{any}
              coupler crossed it. On a long train, that is
              \alert{almost always}.}
      \end{itemize}

      \vspace{0.4em}
      {\footnotesize Estimated cost of a 10,000-scenario corpus:
       \textbf{$\sim$170 CPU-hours}.}
    \end{column}
  \end{columns}
\end{frame}

% --- 55. Why the port was needed (CORRECTED) -----------------------------
\begin{frame}{What the GPU Port Is Required For}
  \vspace{0.5em}
  \begin{columns}[T,onlytextwidth]
    \begin{column}{0.47\textwidth}
      \begin{beamercolorbox}[wd=\linewidth,sep=7pt]{block body}
        {\small\textbf{Required for}}
      \end{beamercolorbox}
      \vspace{0.4em}
      \begin{itemize}\setlength\itemsep{0.4em}
        \item[] {\footnotesize \textbf{Data generation.} Without it the two
              10,000-scenario corpora at $N$ up to 150 do not exist.}
        \item[] {\footnotesize \textbf{The speedup measurement.} The surrogate's
              speedup must be measured against a fast version of the simulator,
              or the comparison is unfair.}
      \end{itemize}
    \end{column}
    \begin{column}{0.47\textwidth}
      \begin{beamercolorbox}[wd=\linewidth,sep=7pt]{block title}
        {\small\textbf{Not required for}}
      \end{beamercolorbox}
      \vspace{0.4em}
      \begin{itemize}\setlength\itemsep{0.4em}
        \item[] {\footnotesize \textbf{The model.} An earlier design evaluated
              the physics equations $F_{\mathrm{phys}}$ in every training step.
              That design was dropped, and \alert{training no longer evaluates
              them}.}
      \end{itemize}
    \end{column}
  \end{columns}

  \vspace{0.9em}
  \begin{minipage}{0.94\textwidth}\footnotesize
    The port also replaced the adaptive solver with fixed-step RK4. Every run
    now takes $\lceil \text{duration}/\mathrm{d}t \rceil$ steps \emph{regardless
    of the driving regime}, so a run can no longer stall and no timeout is
    needed.
  \end{minipage}
\end{frame}

% --- 56. Throughput -------------------------------------------------------
\begin{frame}{Run Time Does Not Grow with Train Length}
  \vspace{0.3em}
  \begin{columns}[T,onlytextwidth]
    \begin{column}{0.48\textwidth}
      \throughputConsist
    \end{column}
    \begin{column}{0.48\textwidth}
      \throughputBatch
    \end{column}
  \end{columns}

  \vspace{0.6em}
  \begin{minipage}{0.94\textwidth}\small
    \alert{Going from 11 to 200 vehicles has no measurable effect on run time.}
    Run time is set by the number of integration steps, each launched as one
    pre-recorded GPU operation, not by the amount of arithmetic. The GPU is far
    from fully used even at $B = 256$ scenarios and $N = 200$ vehicles.

    \vspace{0.4em}
    \textbf{Consequence for the thesis:} the long-train test set
    (\texttt{ood\_size}, $N = 150$) can now be generated, so the claim that the
    model works for any train length can be tested.
  \end{minipage}
\end{frame}

% --- 57. Precision --------------------------------------------------------
\begin{frame}{Numerical Precision, Measured}
  \vspace{1.0em}
  \begin{columns}[c,onlytextwidth]
    \begin{column}{0.31\textwidth}\centering
      \stattile{1.20$\times$}{Cost of double over single precision. The plan
        assumed $\sim$64$\times$}
    \end{column}
    \begin{column}{0.31\textwidth}\centering
      \stattile{60\%}{Of the 0.05 m/s velocity tolerance consumed by float32
        rounding alone}
    \end{column}
    \begin{column}{0.31\textwidth}\centering
      \stattile{f64 $\to$ f32}{Integrate in double, store in single}
    \end{column}
  \end{columns}

  \vspace{1.6em}
  \begin{minipage}{0.92\textwidth}\small
    Run time is limited by \textbf{per-step overhead, not arithmetic}, so double
    precision costs little here. The original plan assumed single precision;
    measuring the cost reversed that choice.
  \end{minipage}
\end{frame}

% --- 58. Two corpora ------------------------------------------------------
\begin{frame}{Two Corpora}
  \vspace{0.4em}
  \begin{columns}[T,onlytextwidth]
    \begin{column}{0.46\textwidth}
      {\footnotesize
      \begin{tabular}{@{}lll@{}}
        \toprule
         & \texttt{data/v1} & \texttt{data/v2} \\
        \midrule
        Curvature   & off        & \texttt{linear}, $k = 1.0$ \\
        Scenarios   & 10,000     & 10,000 \\
        On disk     & 12.4 GB    & 12.4 GB \\
        Integration & 1336 s     & 1389 s \\
        \bottomrule
      \end{tabular}}

      \vspace{0.8em}
      {\footnotesize \texttt{v2} is the dataset used for training. \texttt{v1} is
      kept so the effect of the curvature term can be \textbf{measured} on the
      full dataset.}
    \end{column}
    \begin{column}{0.50\textwidth}
      \begin{center}
        \stattile{753 million}{Vehicle-Steps per Corpus}
        \vspace{0.4em}
        \stattile{$\sim$25 min}{Wall Clock, Against a $\sim$170 CPU-Hour Estimate}
      \end{center}

      \vspace{0.5em}
      {\footnotesize $N$ 10--150 \textbullet\ $T$ 721--2,397 \textbullet\
       duration 180--599 s \textbullet\ 500 consists \textbullet\ 11 regimes
       \textbullet\ 5 corridors.}
    \end{column}
  \end{columns}
\end{frame}

% --- 59. Schema decisions -------------------------------------------------
\begin{frame}{Two Storage Decisions}
  \vspace{0.2em}
  \begin{columns}[c,onlytextwidth]
    \begin{column}{0.53\textwidth}
      \centering
      \scalebox{0.74}{\diskLayout}
    \end{column}
    \begin{column}{0.44\textwidth}
      \small
      \textbf{Static and dynamic separated.} \texttt{H\_hist} would be
      $[T,N,11]$, but seven channels never change, so they are stored once
      rather than at every step.

      \vspace{0.8em}
      \textbf{Routes referenced, not duplicated.} A corridor is reused across
      ${\sim}2{,}000$ scenarios.

      \vspace{0.8em}
      \alert{\texttt{edge\_dynamic} is stored} even though it can be computed
      from the node state, so it can be used as a training target. Act~VII
      shows why this matters.
    \end{column}
  \end{columns}
\end{frame}

% --- 60. Splits -----------------------------------------------------------
\begin{frame}{Splits}
  \vspace{0.4em}
  \centering
  {\small
  \begin{tabular}{@{}lrl@{}}
    \toprule
    \textbf{Split} & \textbf{$n$} & \textbf{Basis} \\
    \midrule
    \texttt{train}         & 4,043 & corridors line0 / line3 / line4 \\
    \texttt{val}           &   517 & random 10\% of eligible \\
    \texttt{test\_id}      &   496 & random 10\% of eligible \\
    \texttt{ood\_size}     &   955 & consists flagged \texttt{ood\_size} \\
    \texttt{ood\_grade}    & 1,999 & \texttt{route\_line2}, steepest \\
    \texttt{ood\_corridor} & 1,990 & \texttt{route\_line5}, wholly held out \\
    \midrule
    \texttt{control\_eval} & \textcolor{uniCrimson}{\textbf{0}}
                           & \textcolor{uniCrimson}{not buildable; see the caveats} \\
    \bottomrule
  \end{tabular}}

  \vspace{1.0em}
  \begin{minipage}{0.92\textwidth}\small
    Corridor reservations are checked \textbf{before} the consist-level
    \texttt{ood\_size} flag, so a reserved corridor cannot leak into training.
  \end{minipage}
\end{frame}

% --- 61. Validation -------------------------------------------------------
\begin{frame}{Validation, and One Finding That Changed Training}
  \vspace{0.2em}
  \begin{columns}[T,onlytextwidth]
    \begin{column}{0.47\textwidth}
      \vspace{0.5em}
      \scalebox{0.86}{\fdConvergenceChart}
    \end{column}
    \begin{column}{0.49\textwidth}
      \vspace{0.2em}
      \small
      \textbf{Both datasets pass} three kinds of check. The most important:
      the equations are re-implemented separately, using \emph{only} the
      scenario file and its route file, so \alert{a bug in the original code
      is not repeated in the check}.

      \vspace{0.7em}
      Each halving of the step size cuts the residual by ${\sim}4\times$, which
      is second-order convergence. The residual therefore comes from the
      finite-difference estimate, \textbf{not} from the stored trajectory.
    \end{column}
  \end{columns}

  \vspace{0.4em}
  \begin{minipage}{0.96\textwidth}\footnotesize
    \textbf{Consequence for training:} the rate of change
    $\mathrm{d}s/\mathrm{d}t$ should not be estimated from differences between
    stored states, as the schema originally suggested. At 0.25\,s spacing that
    error is \alert{as large as the signal the model must learn}. It is
    recomputed from the equations instead.
  \end{minipage}
\end{frame}

% --- 62. Caveats ----------------------------------------------------------
\begin{frame}{Known Data-Quality Caveats}
  \vspace{0.3em}
  \begin{columns}[T,onlytextwidth]
    \begin{column}{0.44\textwidth}
      {\small\textbf{Overspeed, \texttt{data/v2}}}
      \vspace{0.4em}

      {\footnotesize
      \begin{tabular}{@{}lr@{}}
        \toprule
        Never over the limit & 72.9\% \\
        Over by $>$ 5 m/s    & 17.8\% \\
        Over by $>$ 20 m/s   & 2.4\% \\
        Worst                & 67.8 m/s \\
        \bottomrule
      \end{tabular}}

      \vspace{0.7em}
      {\footnotesize Related: \textbf{22\%} of scenarios roll backward faster
       than 1\,m/s, for the same reason.}
    \end{column}
    \begin{column}{0.52\textwidth}
      \small
      The simulator \textbf{intentionally does not enforce} the speed limit
      (\texttt{route\_vmax}). The limit is information for the controller, not a
      physical force. A coasting train on a long downgrade can keep speeding up
      for the full 600 s.

      \vspace{0.7em}
      Overspeed follows the grade data: 39\% of \texttt{route\_line2} scenarios
      against 2\% of \texttt{route\_line5}. \alert{The physics is correct but the
      input was flawed}; Act~V describes the fix.

      \vspace{0.7em}
      \textbf{These runs are kept.} Whether they are bad data or useful hard
      cases is decided at training time. Keeping only runs within 5 m/s of the
      limit retains ${\sim}81\%$.
    \end{column}
  \end{columns}
\end{frame}
